\section{Method}
\label{sec:method}

\subsection{Setup and Notation}

We train an encoder $f_\theta : \mathcal{X} \to \R^d$ that produces feature vectors $\vect{z} = f_\theta(x)$ without final $\ell_2$ normalization. For each input $x$ two augmented views $v_1, v_2$ are generated and passed through the encoder to produce features $\vect{s}_1, \vect{s}_2 \in \R^d$. Stacking a batch of $n$ samples and both views gives feature matrices $\mat{Z}_1, \mat{Z}_2 \in \R^{n \times d}$. Let $\bar{\vect{z}}_k = \tfrac{1}{n}\sum_i \mat{Z}_k[i, :]$ denote the per-view mean, $\tilde{\mat{Z}}_k = \mat{Z}_k - \mathbf{1}\bar{\vect{z}}_k^\top$ the centered batch, and
\begin{equation}
    \mat{C}_k = \tfrac{1}{n}\tilde{\mat{Z}}_k^\top \tilde{\mat{Z}}_k \in \R^{d \times d}
\end{equation}
the per-view empirical covariance. For numerical stability in low-batch or high-dimensional regimes we use the shrinkage-regularized form $\mat{C}_k^{(\rho)} = (1-\rho)\mat{C}_k + \rho\mat{I}_d$ with a small $\rho > 0$ (default $\rho = 0.01$).

\subsection{Derivation from Gaussian Matching}
\label{sec:derivation}

We want the batch of encoder outputs to look, in its second-order moments, like an isotropic standard Gaussian. The natural discrepancy between two distributions is the Kullback--Leibler divergence. For multivariate Gaussians $\mathcal{N}(\boldsymbol{\mu}_1, \boldsymbol{\Sigma}_1)$ and $\mathcal{N}(\boldsymbol{\mu}_2, \boldsymbol{\Sigma}_2)$ on $\R^d$, the closed form is
\begin{equation}
\label{eq:gaussian_kl}
    \mathrm{KL}\!\left(\mathcal{N}(\boldsymbol{\mu}_1, \boldsymbol{\Sigma}_1) \,\big\|\, \mathcal{N}(\boldsymbol{\mu}_2, \boldsymbol{\Sigma}_2)\right)
    \;=\; \tfrac{1}{2}\Bigl[
        \tr\!\bigl(\boldsymbol{\Sigma}_2^{-1}\boldsymbol{\Sigma}_1\bigr)
        + (\boldsymbol{\mu}_2 - \boldsymbol{\mu}_1)^\top \boldsymbol{\Sigma}_2^{-1}(\boldsymbol{\mu}_2 - \boldsymbol{\mu}_1)
        - d
        + \log\!\frac{\det\boldsymbol{\Sigma}_2}{\det\boldsymbol{\Sigma}_1}
    \Bigr].
\end{equation}

\paragraph{Specialization to the isotropic target.}
Our target is the standard Gaussian: $\boldsymbol{\mu}_2 = \vect{0}$, $\boldsymbol{\Sigma}_2 = \mat{I}_d$. The source is the empirical batch Gaussian with moments $\boldsymbol{\mu}_1 = \bar{\vect{z}}$ and $\boldsymbol{\Sigma}_1 = \mat{C}$. Substituting:
\begin{align}
    \mathrm{KL}\!\left(\mathcal{N}(\bar{\vect{z}}, \mat{C}) \,\big\|\, \mathcal{N}(\vect{0}, \mat{I}_d)\right)
    &\;=\; \tfrac{1}{2}\Bigl[\tr(\mat{I}^{-1}\mat{C}) + \bar{\vect{z}}^\top \mat{I}^{-1} \bar{\vect{z}} - d + \log\!\frac{\det\mat{I}}{\det\mat{C}}\Bigr] \notag\\
    \label{eq:full_kl}
    &\;=\; \tfrac{1}{2}\Bigl[\tr(\mat{C}) - \logdet(\mat{C}) - d\Bigr] \;+\; \tfrac{1}{2}\|\bar{\vect{z}}\|^2.
\end{align}
The right-hand side decomposes cleanly into a covariance term---what we call $\mathcal{L}_\text{covKL}$---and a mean term $\tfrac{1}{2}\|\bar{\vect{z}}\|^2$:
\begin{equation}
\label{eq:covkl_derived}
    \mathcal{L}_\text{covKL}(\mat{C}) \;=\; \tfrac{1}{2}\Bigl(\tr(\mat{C}) \;-\; \logdet(\mat{C}) \;-\; d\Bigr),
    \qquad
    \mathcal{L}_\mu(\bar{\vect{z}}) \;=\; \tfrac{1}{2}\|\bar{\vect{z}}\|^2.
\end{equation}
These are precisely the two regularizers of our training loss \eqref{eq:loss}. We separate them because the mean and covariance are structurally different (a vector and a matrix) and admit independent coefficients in practice, but the derivation shows they are not independent design choices: they are the two natural components of one KL to the standard Gaussian.

\paragraph{Connection to the Bregman log-determinant divergence.}
There is a second, equivalent, purely matrix-analytic route to the same object. The function $F(\mat{X}) = -\logdet(\mat{X})$ is smooth and strictly convex on the cone of symmetric positive-definite matrices, with gradient $\nabla F(\mat{X}) = -\mat{X}^{-1}$. Its Bregman divergence (also known as \emph{Stein's loss} or the LogDet divergence) is
\begin{equation}
    D_F(\mat{X}, \mat{Y}) \;=\; F(\mat{X}) - F(\mat{Y}) - \bigl\langle \nabla F(\mat{Y}),\, \mat{X} - \mat{Y}\bigr\rangle
    \;=\; \tr(\mat{Y}^{-1}\mat{X}) \;-\; \log\!\frac{\det\mat{X}}{\det\mat{Y}} \;-\; d.
\end{equation}
Setting $\mat{Y} = \mat{I}_d$ gives $D_F(\mat{X}, \mat{I}) = \tr(\mat{X}) - \logdet(\mat{X}) - d = 2\,\mathcal{L}_\text{covKL}(\mat{X})$. So our covariance objective is (up to a factor of two) the Bregman divergence of the log-determinant potential from the identity to the empirical covariance---a natural geometric object on the cone of covariance matrices, independent of the Gaussian-KL framing.

\paragraph{Term-by-term interpretation.}
Each of the three components of $\mathcal{L}_\text{covKL}$ has a clean geometric meaning:
\begin{itemize}
    \item $\tfrac{1}{2}\tr(\mat{C}) = \tfrac{1}{2}\sum_j \sigma_j^2$ is half the total variance. It is the scale-anchoring term: grows quadratically with feature scale, vanishes at zero.
    \item $-\tfrac{1}{2}\logdet(\mat{C}) = -\sum_i \tfrac{1}{2}\log\lambda_i$ is (up to additive constants) the negative differential entropy of $\mathcal{N}(\vect{0}, \mat{C})$. It is the anti-collapse term: diverges as any eigenvalue approaches zero, rewarding spread across directions.
    \item The constant $-\tfrac{d}{2}$ is the ``baseline''---it normalizes $\mathcal{L}_\text{covKL}(\mat{I}_d) = 0$.
\end{itemize}
The nontrivial work of the objective is in the interplay between the first two: trace pulls scale down, $-\logdet$ pushes spread up, and both are simultaneously optimal at $\mat{C} = \mat{I}_d$.

\subsection{The CovKL Objective}

Having derived it, we now collect its key properties. The covariance-KL objective is
\begin{equation}
\label{eq:covkl_def}
    \mathcal{L}_\text{covKL}(\mat{C}) \;\triangleq\; \tfrac{1}{2}\Bigl(\tr(\mat{C}) \;-\; \logdet(\mat{C}) \;-\; d\Bigr).
\end{equation}
For any symmetric positive-definite $\mat{C}$, this quantity equals the KL divergence $\mathrm{KL}(\mathcal{N}(\vect{0}, \mat{C}) \,\|\, \mathcal{N}(\vect{0}, \mat{I}_d))$. It is nonnegative, zero only at $\mat{C} = \mat{I}_d$, and a smooth function of $\mat{C}$'s eigenvalues:
\begin{equation}
\label{eq:covkl_eigendecomp}
    \mathcal{L}_\text{covKL}(\mat{C}) \;=\; \tfrac{1}{2}\sum_{i=1}^{d}\bigl(\lambda_i - \log\lambda_i - 1\bigr),
\end{equation}
where $\lambda_1 \ge \dots \ge \lambda_d > 0$ are the eigenvalues of $\mat{C}$. Each term $\lambda - \log\lambda - 1$ is convex, vanishes at $\lambda = 1$, and diverges as $\lambda \to 0^+$ (via $-\log\lambda$) and as $\lambda \to \infty$ (via the $\lambda$ term). The objective is therefore a two-sided barrier on every eigenvalue, pulling all of them toward one.

\paragraph{Scale / structure decomposition.}
Writing $\mat{C} = \mat{D}^{1/2}\mat{R}\mat{D}^{1/2}$ with $\mat{D} = \diag(\sigma_1^2, \dots, \sigma_d^2)$ and $\mat{R}$ the correlation matrix,
\begin{equation}
\label{eq:decomp}
    \mathcal{L}_\text{covKL}(\mat{C}) \;=\; \underbrace{\tfrac{1}{2}\sum_{j=1}^{d}\bigl(\sigma_j^2 - 2\log\sigma_j - 1\bigr)}_{\mathcal{L}_\text{diag}} \;\underbrace{-\; \tfrac{1}{2}\logdet(\mat{R})}_{\mathcal{L}_\text{corr}}.
\end{equation}
The diagonal term $\mathcal{L}_\text{diag}$ controls per-dimension scale: it vanishes at $\sigma_j = 1$, diverges as $\sigma_j \to 0$ or $\sigma_j \to \infty$, and is locally quadratic, $\mathcal{L}_\text{diag} \approx \sum_j (\sigma_j - 1)^2$ near the optimum. The correlation term $\mathcal{L}_\text{corr}$ penalizes small eigenvalues of $\mat{R}$: it vanishes at $\mat{R} = \mat{I}$ (subject to $\tr\mat{R} = d$) and diverges as any eigenvalue of $\mat{R}$ approaches zero.

\subsection{Relation to Existing Flat-Space SSL Losses}
\label{sec:approximations}

We now recover the scale and decorrelation terms of VICReg, Barlow Twins, and SimDINO as approximations of $\mathcal{L}_\text{diag}$ and $\mathcal{L}_\text{corr}$.

\paragraph{SigReg vs.\ $\mathcal{L}_\text{diag}$.}
SimDINO's two-sided variance penalty \citep{schaeffer2024sigreg} is $\tfrac{1}{d}\sum_j (\sigma_j - 1)^2$. Taylor-expanding $\sigma_j^2 - 2\log\sigma_j - 1$ around $\sigma_j = 1$:
\begin{equation}
    \sigma_j^2 - 2\log\sigma_j - 1 \;=\; 2(\sigma_j - 1)^2 \;-\; \tfrac{4}{3}(\sigma_j - 1)^3 \;+\; O\!\bigl((\sigma_j - 1)^4\bigr).
\end{equation}
So SigReg agrees with $\mathcal{L}_\text{diag}$ to leading order, up to an overall factor of 2. The disagreement is at the tails: as $\sigma_j \to 0$, SigReg saturates at 1 while $\mathcal{L}_\text{diag}$ diverges. The SigReg penalty provides no barrier against per-dimension collapse; $\mathcal{L}_\text{diag}$ provides an unbounded one.

\paragraph{VICReg hinge vs.\ $\mathcal{L}_\text{diag}$.}
VICReg's hinge $\max(0, 1 - \sigma_j)$ is asymmetric---zero for $\sigma_j \ge 1$, linear below. This differs structurally from $\mathcal{L}_\text{diag}$, which penalizes both tails. The asymmetry reflects a pragmatic design choice: a direction whose variance grows above one is presumed to carry useful signal and should not be suppressed. $\mathcal{L}_\text{diag}$ respects this intuition differently: its upper-side growth is logarithmic ($\sigma_j^2 - 2\log\sigma_j - 1 \sim \sigma_j^2$ for large $\sigma_j$) which, while not zero, is gentler than the quadratic $(\sigma_j - 1)^2$ growth of SigReg. The local effect of $\mathcal{L}_\text{diag}$ upper-side penalty is weaker than SigReg's, addressing the same concern as VICReg's hinge without the one-sided discontinuity.

\paragraph{VICReg covariance vs.\ $\mathcal{L}_\text{corr}$.}
VICReg's decorrelation term is $\tfrac{1}{d}\sum_{i \ne j} \mat{C}_{ij}^2$---the Frobenius norm of the off-diagonal covariance. This is a second-order approximation to $-\logdet\mat{R}$ near $\mat{R} = \mat{I}$: writing $\mat{R} = \mat{I} + \mat{E}$ with $\mat{E}$ small and zero-diagonal,
\begin{equation}
    -\logdet(\mat{I} + \mat{E}) \;=\; \tfrac{1}{2}\|\mat{E}\|_F^2 \;+\; O(\mat{E}^3),
\end{equation}
since the linear term $\tr(\mat{E}) = 0$ vanishes by $\mat{E}$'s zero diagonal. So VICReg's Frobenius off-diagonal is the quadratic approximation to $\mathcal{L}_\text{corr}$. The approximation agrees near $\mat{R} = \mat{I}$ but misses the $-\logdet$ barrier against directional collapse.

\paragraph{Coding rate vs.\ $\mathcal{L}_\text{corr}$.}
SimDINO and MCR\textsuperscript{2}'s coding-rate expansion is $\tfrac{1}{2}\logdet(\mat{I}_d + (d/\varepsilon^2)\mat{R})$. This term \emph{maximizes} $\logdet(\mat{I} + \alpha \mat{R})$ for $\alpha = d/\varepsilon^2$, which rewards spread but does not have a minimum. Specifically, if $\lambda_i(\mat{R}) \to 0$ for some direction, the coding-rate factor $1 + \alpha\lambda_i \to 1$ and the contribution saturates---the gradient resisting near-singular correlation is $O(\alpha)$, large but bounded. The corresponding $-\logdet\mat{R}$ diverges unboundedly. Near $\mat{R} = \mat{I}$ the two terms have similar gradients, but in the collapse regime they behave qualitatively differently.

\subsection{Why Coding Rate Requires a Scale Anchor}
\label{sec:coding_rate_scale}

A concern with using SimDINO-style coding rate on the \emph{raw covariance} $\mat{C}$ (rather than the correlation $\mat{R}$) without $\ell_2$ normalization is that the objective is unbounded as a function of feature scale. We sketch this argument informally; a fully rigorous treatment would be uninteresting.

\paragraph{Scale sensitivity.}
Let $f_\theta$ be an encoder with output features $\vect{z}$, and consider the rescaled encoder $f_\theta^{(\alpha)}(x) = \alpha f_\theta(x)$ with $\alpha > 0$. The batch covariance transforms as $\mat{C}^{(\alpha)} = \alpha^2 \mat{C}$. The coding rate is
\begin{align}
    R_\varepsilon(\alpha^2 \mat{C}) \;=\; \tfrac{1}{2}\logdet\!\Bigl(\mat{I}_d + \tfrac{d}{n\varepsilon^2}\alpha^2 \mat{C}\Bigr)
    \;=\; \tfrac{1}{2}\sum_{i=1}^{d} \log\!\Bigl(1 + \tfrac{d\,\alpha^2}{n\varepsilon^2}\lambda_i(\mat{C})\Bigr),
\end{align}
which, for any fixed full-rank $\mat{C}$, grows as $\alpha \to \infty$:
\begin{equation}
    R_\varepsilon(\alpha^2 \mat{C}) \;\sim\; d\,\log(\alpha) \;+\; \text{const} \;\to\; \infty.
\end{equation}
A SimDINO-style loss of the form $\mathcal{L} = \mathcal{L}_\text{align} - \gamma R_\varepsilon(\mat{C})$ is therefore monotonically improved by scaling up the encoder output, up to whatever counter-pressure the alignment term provides.

\paragraph{What does alignment do against this?}
Cross-view squared-Euclidean alignment $\mathcal{L}_\text{align} = \|\vect{s}_1 - \vect{s}_2\|^2$ scales as $\alpha^2$: if views differ on average by $\Delta$ per dimension at scale $\alpha = 1$, they differ by $\alpha\Delta$ at scale $\alpha$. The total loss then has
\begin{equation}
    \mathcal{L}(\alpha) \;=\; \alpha^2 \mathcal{L}_\text{align}^{(0)} \;-\; \gamma\bigl(d\log\alpha + O(1)\bigr),
\end{equation}
which is convex in $\alpha$ with a minimum at $\alpha^\star = \sqrt{\gamma d / (2\mathcal{L}_\text{align}^{(0)})}$. So alignment \emph{does} eventually win at large scale, but the equilibrium scale depends on the ratio $\gamma d / \mathcal{L}_\text{align}^{(0)}$: the loss has no intrinsic scale, and training dynamics determine where in scale space the encoder lives. In practice this manifests as sensitivity to $\gamma$ and gradual drift of $\|\vect{z}\|$ over training.

\paragraph{Contrast with $\mathcal{L}_\text{covKL}$.}
For the covariance KL, applying the same rescaling:
\begin{equation}
    \mathcal{L}_\text{covKL}(\alpha^2\mat{C}) \;=\; \tfrac{1}{2}\Bigl(\alpha^2 \tr(\mat{C}) \;-\; \logdet(\mat{C}) \;-\; 2d\log\alpha \;-\; d\Bigr),
\end{equation}
which diverges to $+\infty$ as $\alpha \to \infty$ (dominated by the $\alpha^2 \tr(\mat{C})$ term) \emph{and} as $\alpha \to 0^+$ (dominated by $-2d\log\alpha$). The $\tr(\mat{C})$ term is precisely the scale anchor that coding rate lacks: scaling up features is now expensive, not free. The equilibrium $\alpha^\star = 1$ is determined by the objective itself, not by a balance against a separate alignment term. This is the practical consequence of $\mathcal{L}_\text{covKL}$ being a \emph{matching} rather than \emph{expansion} objective.

\paragraph{Practical implication.}
This is why SimDINO, when used without $\ell_2$ normalization, requires a separate scale regularizer (SigReg or a correlation-rate variant that standardizes $\mat{C}$ to $\mat{R}$). The coding-rate term's scale unboundedness is not a bug per se---paired with proper scale control, it trains successfully---but it does mean the two-term recipe is load-bearing: removing either collapses training. $\mathcal{L}_\text{covKL}$ replaces both terms with one that is self-stabilizing in scale.

\paragraph{Barlow Twins.}
Barlow's cross-correlation objective is not directly in the form of \eqref{eq:decomp}: it operates on the cross-view object $\mat{C}^\text{xy} = \mat{Z}_1^\text{bn}\!^\top\mat{Z}_2^\text{bn}/n$, not a within-view covariance. But Barlow's on-diagonal $(1 - \mat{C}^\text{xy}_{jj})^2$ fuses invariance and scale control (the diagonal of $\mat{C}^\text{xy}$ is the per-dimension cross-view cosine, which equals one only if both features are unit-scale \emph{and} aligned), while its off-diagonal ${\mat{C}^\text{xy}_{ij}}^2$ simultaneously decorrelates dimensions and aligns cross-view. So Barlow's single term does the work of VICReg's three. We view this as an alternative unification, dual to ours: Barlow merges invariance into the regularizer, while CovKL keeps invariance separate and merges the two regularizer terms.

\subsection{Taylor Approximations and Computational Variants}
\label{sec:taylor_variants}

The eigenvalue form \eqref{eq:covkl_eigendecomp} invites systematic approximation by Taylor expansion around $\lambda = 1$. We develop the expansion here and show that the leading-order form reconstructs, exactly, a natural cousin of VICReg's covariance objective---revealing VICReg's decorrelation term as an approximate Gaussian-matching penalty. We also discuss why these approximations are worth treating as first-class training objectives rather than mere conceptual curiosities: their computational cost can be substantially lower than that of $\mathcal{L}_\text{covKL}$ itself.

\paragraph{Second-order Taylor expansion.}
The per-eigenvalue summand is
\begin{equation}
    f(\lambda) \;=\; \lambda - \log\lambda - 1.
\end{equation}
A Taylor expansion around $\lambda = 1$ gives
\begin{equation}
\label{eq:taylor}
    f(\lambda) \;=\; \tfrac{1}{2}(\lambda - 1)^2 \;-\; \tfrac{1}{3}(\lambda - 1)^3 \;+\; \tfrac{1}{4}(\lambda - 1)^4 \;-\; \cdots,
\end{equation}
obtained from $f(1) = 0$, $f'(1) = 0$, and $f^{(k)}(1) = (-1)^k(k-1)!$ for $k \ge 2$. Truncating at second order and summing over eigenvalues,
\begin{equation}
\label{eq:covkl_f_approx}
    \mathcal{L}_\text{covKL}(\mat{C}) \;\approx\; \tfrac{1}{4} \sum_{i=1}^{d} (\lambda_i - 1)^2 \;=\; \tfrac{1}{4} \|\mat{C} - \mat{I}_d\|_F^2,
\end{equation}
where the last equality uses the orthogonal invariance of the Frobenius norm: $\|\mat{C} - \mat{I}\|_F^2 = \sum_i (\lambda_i - 1)^2$ for any symmetric $\mat{C}$. So the second-order approximation of $\mathcal{L}_\text{covKL}$ is simply a quarter of the squared Frobenius distance from the empirical covariance to the identity.

\paragraph{Diagonal--off-diagonal split.}
Expanding the Frobenius norm componentwise makes the approximation's structure visible:
\begin{equation}
\label{eq:frob_split}
    \|\mat{C} - \mat{I}_d\|_F^2 \;=\; \underbrace{\sum_{j=1}^{d}\bigl(\sigma_j^2 - 1\bigr)^2}_{\text{diagonal: scale}} \;+\; \underbrace{\sum_{i \ne j} \mat{C}_{ij}^2}_{\text{off-diagonal: decorrelation}}.
\end{equation}
The off-diagonal term is exactly VICReg's decorrelation loss \citep{bardes2022vicreg}. The diagonal term is a scale penalty of the form $(\sigma_j^2 - 1)^2$---which differs from both SimDINO's SigReg (which uses $(\sigma_j - 1)^2$) and from VICReg's own asymmetric hinge, but sits in the same family. We can read \eqref{eq:frob_split} as \textbf{the VICReg decorrelation loss paired with its ``Gaussian-matching'' natural diagonal companion}, derived as a single second-order Taylor approximation of $\mathcal{L}_\text{covKL}$. This is a stronger justification for VICReg's covariance objective than the original paper offers: it is the quadratic piece of an information-theoretic divergence, not an ad-hoc redundancy penalty.

\paragraph{Approximation family.}
Each truncation order of \eqref{eq:taylor} defines a candidate training objective. We name them for ease of reference:
\begin{description}
    \item[\textsc{CovKL-F} (Frobenius, order~2):] $\tfrac{1}{4}\|\mat{C} - \mat{I}_d\|_F^2$. The leading-order approximation. Computable with only matrix multiplies; no eigendecomposition, no $\logdet$.
    \item[\textsc{CovKL-F3} (order~3):] $\tfrac{1}{4}\|\mat{C} - \mat{I}_d\|_F^2 - \tfrac{1}{6}\sum_i(\lambda_i - 1)^3$. The cubic correction requires the eigendecomposition, giving up the computational advantage.
    \item[\textsc{CovKL} (exact):] $\tfrac{1}{2}(\tr\mat{C} - \logdet\mat{C} - d)$. Requires a Cholesky or LU decomposition for $\logdet$.
    \item[\textsc{CovKL-Diag} (diagonal-only, order~2):] $\tfrac{1}{4}\sum_j(\sigma_j^2 - 1)^2$. The scale piece of CovKL-F with no off-diagonal penalty. Computable in $O(d)$. Useful as an ablation to isolate the scale contribution.
\end{description}

\paragraph{Computational cost.}
For a batch of $n$ samples in $\R^d$ with $n \ge d$:
\begin{itemize}
    \item Computing $\mat{C}$: $O(n d^2)$, unavoidable for all variants.
    \item $\tr(\mat{C})$: $O(d)$.
    \item $\|\mat{C} - \mat{I}_d\|_F^2$: $O(d^2)$, elementwise; no decomposition.
    \item $\logdet(\mat{C})$: $O(d^3)$ via Cholesky, plus differentiability through PyTorch's \texttt{slogdet} which is moderately heavy in practice and can be numerically unstable near singular $\mat{C}$ (whence the shrinkage $\rho$).
    \item Eigenvalues $\{\lambda_i\}$: $O(d^3)$ via symmetric eigendecomposition.
\end{itemize}
At $d = 256$, the nominal gap between $O(d^2)$ and $O(d^3)$ is roughly $250\times$; in practice the gap is smaller because highly optimized BLAS routines are closer in wall-clock time per FLOP than the asymptotic counts suggest. But for larger projection-head dimensions $d \ge 1024$, the cubic cost of $\logdet$ becomes a meaningful fraction of the training step, and the Frobenius approximation is attractive.

\paragraph{Accuracy of the approximation.}
The second-order Frobenius approximation and the exact CovKL agree to leading order near $\mat{C} = \mat{I}$, but diverge in two ways that matter at training time. First, CovKL-F is bounded as any $\lambda_i \to 0$, while the exact CovKL diverges. The Frobenius approximation therefore provides a weaker barrier against dimensional collapse---the same tradeoff that distinguishes VICReg's off-diagonal covariance loss from a full $-\logdet\mat{R}$ barrier in \secref{sec:approximations}. Second, CovKL-F grows as $\lambda^2$ for large $\lambda$, while the exact CovKL grows as $\lambda - \log\lambda \sim \lambda$---so the Frobenius form is more aggressive on over-scaled directions. In our experiments (\secref{sec:method} and forthcoming), CovKL-F trains stably and closely tracks the exact CovKL on CIFAR-10, at a fraction of the compute. At larger scales and with longer training, the collapse-barrier difference may become more important, which we flag as a question for future empirical work.

\subsection{Training Procedure}

The training loss combines a standard invariance term, a mean penalty, and $\mathcal{L}_\text{covKL}$, applied per view:
\begin{equation}
\label{eq:loss}
    \mathcal{L} \;=\; \lambda_\text{align}\,\tfrac{1}{2}\bigl(\|\vect{s}_1 - \vect{s}_2\|^2\bigr) \;+\; \tfrac{\lambda_\mu}{2 d}\bigl(\|\bar{\vect{z}}_1\|^2 + \|\bar{\vect{z}}_2\|^2\bigr) \;+\; \tfrac{\lambda_C}{2 d}\bigl(\hat{\mathcal{L}}_\text{covKL}(\mat{C}_1^{(\rho)}) + \hat{\mathcal{L}}_\text{covKL}(\mat{C}_2^{(\rho)})\bigr),
\end{equation}
where $\hat{\mathcal{L}}_\text{covKL}$ is the unnormalized form in \eqref{eq:covkl_def} and we divide by $d$ so each term is $O(1)$ in feature dimension. Default coefficients are $\lambda_\text{align} = 25$, $\lambda_\mu = 1$, $\lambda_C = 1$, $\rho = 10^{-2}$; these were not tuned beyond an initial scale-matching check at random initialization. No EMA teacher and no stop-gradient are used---the model is a plain Siamese network with weight sharing.

\paragraph{Per-view rather than pooled.}
We compute $\mat{C}_k^{(\rho)}$ and $\bar{\vect{z}}_k$ separately per view rather than on the concatenation $\mat{Z} = [\mat{Z}_1; \mat{Z}_2]$. Pooling doubles the sample count for the covariance estimate but introduces artificial within-sample correlation: the same image contributes twice through correlated views. Per-view estimates are unbiased; pooled estimates are not. In practice the difference is small for large $n$, but the per-view form is cleaner.

\paragraph{Shrinkage $\rho$.}
At random initialization the empirical covariance is far from $\mat{I}$ and the $-\logdet$ term can be numerically unstable. Shrinkage $\mat{C}^{(\rho)} = (1-\rho)\mat{C} + \rho\mat{I}$ bounds the smallest eigenvalue at $\rho\cdot(\text{min eigenvalue of } \mat{C})$ plus the shrinkage floor, giving a finite $\logdet$ at initialization and a smooth training trajectory. $\rho = 10^{-2}$ is sufficient in practice; smaller $\rho$ converge to the same solution but warm up more slowly.

\paragraph{Dimension-normalization of regularizers.}
Both $\mathcal{L}_\text{covKL}$ (scales as $O(d)$ from the $d$ eigenvalue contributions) and $\|\bar{\vect{z}}\|^2$ (scales as $O(d)$ if per-dimension means are $O(1)$) are divided by $d$ to yield $O(1)$-scale terms. Without this normalization, $\mathcal{L}_\text{covKL}$ dominates the loss at random initialization (where $\sim d$ eigenvalues are far from one) and the invariance term has negligible gradient. Empirically, skipping the $d$-normalization causes early collapse of the student representation.

\subsection{Summary of Derivation}

We have shown that existing flat-space SSL recipes can be read as two-term approximations of a single Gaussian second-order KL objective:
\begin{itemize}
    \item SigReg's two-sided $(\sigma_j - 1)^2$ is the leading Taylor term of $\sigma_j^2 - 2\log\sigma_j - 1$ near $\sigma_j = 1$.
    \item VICReg's off-diagonal Frobenius covariance is the leading Taylor term of $-\logdet\mat{R}$ near $\mat{R} = \mat{I}$.
    \item MCR\textsuperscript{2} coding-rate $\tfrac{1}{2}\logdet(\mat{I} + (d/\varepsilon^2)\mat{R})$ is a structurally different object---an \emph{expansion} rather than a matching objective---and has no minimum without an additional scale anchor.
\end{itemize}
The unified object $\mathcal{L}_\text{covKL}$ combines the advantages: a true minimum at $\mat{C} = \mat{I}$, two-sided scale barriers on each eigenvalue, no separate scale regularizer, and a gentler (logarithmic) upper-side penalty on per-dimension variance than SigReg's quadratic. The training recipe in \eqref{eq:loss} replaces the two-or-three term regularizer of each existing method with the single $\mathcal{L}_\text{covKL}$ term.
